\cotefile{45}
\scanpage{2}
Groupes \unsure{semi-simples|semi simples}
% archive sheet number: 1

\scanpage{3}
\unsure{Application} des \unsure{ss-Alg.|ss.-Alg.|s.-Alg.} de \unsure{Cartan} aux corps de dim 1.
% archive sheet number: 2

\scanpage{4}
\uline{Lemme 1} Soit $k$ un corps, et supposons que les conditions suivantes soient \unsure{vérifiées|vérifiée}

\deleted{\unsure{1)}} On a $H^2(k,\unsure{\mathscr{K}|\mathcal{K}})=0$ dans les cas suivants \inserted{($\mathscr{K}$ commutatif)}

\begin{enumerate}[label=\alph*)]
\item $\mathscr{K}$ un tore
\item $\mathscr{K}$ \deleted{un \unsure{schéma}} fini ét. sur $k$
\item $\mathscr{K}$ une V.A.
\item $\mathscr{K}$ groupe unipotent connexe annulé par $p$.
\end{enumerate}

Alors pour tte \deleted{\unsure{donnée}} \inserted{suite exacte} de \inserted{\unsure{schémas en}} groupes Alg.
\diagram{\begin{tikzcd}[ampersand replacement=\&]
e \arrow[r,dash] \& K \arrow[r,dash] \& \unsure{G|C} \arrow[r,dash] \& G' \arrow[r,dash] \& e
\end{tikzcd}}
(\unsure{non} commutatifs), tels que \deleted{$K$} \inserted{$G,G'$} \uline{simples/$k$}, on a
\[
H^1(k,G)\longrightarrow H^1(k,G')\qquad\text{\uline{surjectif}}
\]
% slanted
\margin{Inutile de \unsure{supposer} $G,G'$ \unsure{de t.f.|simples}. On se \unsure{ramène} \unsure{d'abord} \uline{\unsure{au}} \uline{\unsure{cas}} \uline{$K$ \unsure{réduit}} \unsure{? Si oui}}

Principe de démonstration. On prend un tore maximal $T$ de $K$ défini sur $k$, et utilisant \unsure{le|les} th. de conjugaison, on voit que $N_G(T)\longrightarrow G'$ est \deleted{surjectif} \inserted{\unsure{épim.}}, donc on aura
% slanted
\margin{faut-il supposer $K$ linéaire??? On \unsure{peut} commencer par \unsure{dévisser} en \uline{groupe fini}, et \deleted{groupe} \inserted{\unsure{groupes}} \uline{commutatifs} \uline{connexes} (\unsure{ou commutables}) \uline{et} groupe \uline{\unsure{quasi-réductif}} \unsure{Opérer ici leur propre} \unsure{comm.} \unsure{aux} \uline{\unsure{linéaires}}}
\diagram{\begin{tikzcd}[ampersand replacement=\&]
e \arrow[r,dash] \& N_K(T) \arrow[r,dash] \& N_G(T) \arrow[r,dash] \& G' \arrow[r,dash] \& e
\end{tikzcd}}
donc on peut supposer \uline{$T$ invariant dans $G$}, a fortiori dans $K$. Car on \unsure{se} ramène aux deux cas

1) $K$ est un tore \deleted{\unsure{1)}} $K$ est \deleted{sans \unsure{tore}} \inserted{de \unsure{rang}} \unsure{réductif} 0. \unsure{Par} dévissage ce dernier cas se réduit à :

2) $K$ est fini ét. /$k$ \deleted{\unsure{3)}} 3) $K$ est connexe, et géométriquement une ext. d'une VA par un unipotent. Par dévissage avec les commutateurs successifs, on est ramené
% archive sheet number: 3

\scanpage{5}
\inserted{\unsure{dans 1°}} au cas d'un groupe commutatif, \deleted{\unsure{p}} et par dévissage encore au cas où ce groupe \deleted{\unsure{vérifie}} soit $pK=K$, soit $pK=0$, i.e. les cas 3) $K$ une \unsure{VA.|V.A.}

4) $K$ est unipotent comm. connexe annulé par $p$.

Pour le cas 2), i.e. $K$ fini étale sur $k$, pas néc. commutatif, on note que l'on a une suite exacte (où $\bar{k}$ est la clôture \deleted{Algébrique} séparable de $k$)
\diagram{\begin{tikzcd}[ampersand replacement=\&]
e \arrow[r] \& K(\bar{k}) \arrow[r] \& \unsure{G|C}(\bar{k}) \arrow[r,dash] \& G'(\bar{k}) \arrow[r,dash] \& e
\end{tikzcd}}
et en vertu de l'argument de \unsure{Schreier|Schur}--Serre, on aura
\[
H^1(\pi,G(\bar{k}))-H^1(\pi,G'(\bar{k}))\quad\text{surjectif si}
\]
$H^2(\pi,A)=0$ pour tt $\pi$-module fini $A$. Or cette hypothèse signifie $H^2(k,K)$ pour tt $K$ fini étale sur $k$, et pour un groupe $G$ simple sur $k$, on a
\[
\uline{H^i(k,G)\simeq H^i(\pi,G(\bar{k}))},
\]
\inserted{pour tt $i$, en particulier \unsure{lorsque} $i=\unsure{1|0}$} ce qui prouve le lemme.

\uline{Corollaire} Supposons que $H^2(k,K)=0$ pour tt $K$ fini sur $k$ (pas néc. étale /$k$). Alors les \deleted{conditions} hyp du lemme 1 sont vérifiées, donc aussi la conclusion

Notons d'abord que si $G$ simple sur $k$, alors les $H^i(k,G)$ sont torsion pour $i\unsure{\geq|>}0$ [car $=H^i(\pi,G(\bar{k}))$] donc \deleted{\unsure{si $G$ les|si $G$ est}} \inserted{\unsure{tt}} \unsure{\{|f}\unsure{c}
% archive sheet number: 4

\scanpage{6}
$H^i(k,\unsure{G|C})$ provient d'un $H^i(k,{}_{\unsure{n|-}}G)$. Or si $G$ est une VA ou un tore, ${}_nG$ est fini sur $k$. Cela montre que les conditions du lemme 1 sont vérifiées pour a) b) c). Reste à vérifier d).

Pour \unsure{ceci|cela}, on va \deleted{\unsure{supposer}} utiliser \deleted{\unsure{un}} le fait que $H^2(k,U)$ devient nul sur une extension de $k$ de la forme $k^{p^{-n}}$, [car $U$ \unsure{splitte} sur une telle extension en groupes additifs] et on va prouver le

\uline{Lemme 2} Pour tt schéma en groupes $G$ sur $k$, le noyau de $H^i(k,G)-H^i(k^{p^{-n}},G)$ est égal à celui de
\[
H^i(k,G)\xrightarrow{F^n}H^i(k,G^{(p^n)})
\]
défini par $F^n$ ($F=\mathrm{Frob}$). Plus précisément, on a commutativité dans le diagramme
% slanted
\margin{\unsure{Pourrait} s'utiliser directement, $+$ \unsure{cette} \unsure{déduction} \unsure{de} \unsure{truc} \unsure{pourrait} \unsure{ramener} au cas du groupe $K$. \deleted{\unsure{faudrait}} vérifier ???}

\deleted{[N.B.}
\diagram{\begin{tikzcd}[ampersand replacement=\&]
H^i(k,G) \arrow[r] \arrow[dr,"F^n"'] \& H^i(k^{p^{-n}},G)\deleted{\simeq H^i(k,G^{(p^n)})} \arrow[d,dash,"\wr"] \\
\& H^i(k,G^{(p^n)})
\end{tikzcd}}
\diagram{\begin{tikzcd}[ampersand replacement=\&,column sep=small]
G \arrow[d,"p"'] \& G^{(p^n)} \arrow[l,"p^n"'] \arrow[d] \arrow[dl,phantom,"\text{cart.}" description] \& G \arrow[l,"F^n"'] \arrow[ll,bend right=38,"\unsure{\varphi_G|\phi_G}"'] \arrow[dl,"p"] \\
\begin{gathered}S\\ \Vert\\ \operatorname{Sp}(k)\end{gathered} \& \begin{gathered}S\\ \Vert\\ \operatorname{Sp}(k)\end{gathered} \arrow[l,"\varphi_S"] \&
\end{tikzcd}}
\[
\lambda\longmapsto\lambda^{p^n}
\]
Pour le voir, on est ramené au cas du $H^0$, où c'est trivial [N.B. On n'a pas besoin que $G$ soit \unsure{représentable}\ldots].
% archive sheet number: 5

\scanpage{7}
\uline{Corollaire} Ledit noyau est donc égal à l'image de $H^i(k,\operatorname{Ker}(F_{\unsure{G|C}}^n))$, $F_G^n:G-G^{(p^n)}$ le Frobenius itéré.

Or si $G$ est simple /$k$, alors $F_G^n$ est radiciel, donc son noyau est radiciel /$k$, a fortiori fini. Cela achève de prouver le Corollaire au lemme 1.

\uline{Lemme 3} Pour que $H^1(k,G)-H^1(k,G')$ soit un épim pour tt épim $G\xrightarrow{u}G'$ de schémas en groupes algébriques, il suffit qu'il en soit ainsi dans les deux cas suivants

\begin{enumerate}[label=\alph*)]
\item $G$, \deleted{$\unsure{K}$} et $K=\operatorname{Ker}u$ \deleted{\unsure{sont}} \inserted{(comme $G'$)} simples /$k$
\item $K$ est radiciel.
\end{enumerate}

Dans le cas général, il existe en effet un sous groupe radiciel \inserted{inv. $I$} de $G$ \unsure{[|(}noyau d'un Frobenius itéré convenable] tel que $G/I$ \deleted{$\unsure{K}$} et $K/I\cap K$ sont simples /$k$.

Considérons
\diagram{\begin{tikzcd}[ampersand replacement=\&]
e \arrow[r] \& K \arrow[r,dash] \arrow[d] \& G \arrow[r,dash,"\unsure{f|j}"] \arrow[d,"u"'] \& G/K \arrow[r] \arrow[d,"v"] \& e \\
e \arrow[r,dash] \& K/K\cap I \arrow[r,dash] \& G/I \arrow[r,dash,"f'"] \& G/KI \arrow[r,dash] \& e
\end{tikzcd}}
\deleted{Alors} d'où
% archive sheet number: 6

\scanpage{8}
\diagram{\begin{tikzcd}[ampersand replacement=\&]
H^1(k,I) \arrow[r,dash] \arrow[d] \& H^1(k,I/I\cap K) \arrow[d] \\
H^1(k,\unsure{G|C}) \arrow[r,dash,"\unsure{f|j}_1"] \arrow[d,"u_1"'] \& H^1(k,G/K) \arrow[d,"v_1"] \\
H^1(k,G/I) \arrow[r,dash,"f'_1"] \& H^1(k,G/KI)
\end{tikzcd}}
Par l'hypothèse \inserted{a)} $u_1$ et $f'_1$ sont surjectifs, donc $f'_1u_1=v_1f_1$ l'est, donc si $\xi'\in H^1(k,G/K)$, il existe $\xi\in H^1(k,G)$ tel que $\xi'$ et $f_1(\xi)$ aient même image dans $H^1(k,G/KI)$. Quitte à tordre par $\xi'$, on peut supposer que $\xi'=0$, donc que $v_1(\xi)=0$, donc que $\xi$ est \deleted{dans} l'image d'un $\unsure{\chi|\zeta}\in H^1(k,KI/K)=H^1(k,I/I\cap K)$.

Or par hypothèse b), \deleted{\unsure{il s'ensuit}} $\chi$ est dans l'image de $H^1(k,I)$, et on gagne.

\uline{Lemme 4} Supposons qu'on ait \deleted{\unsure{a)}}

\unsure{c)|e)} $H^2(k,\unsure{\mathscr{K}|\mathcal{K}})=0$ si $K$ groupe \inserted{(commutatif)} fini radiciel sur $k$.

Alors pour tte \unsure{suite|schémas} exacte de groupes algébriques
\diagram{\begin{tikzcd}[ampersand replacement=\&]
e \arrow[r,dash] \& K \arrow[r,dash] \& G \arrow[r,dash] \& G' \arrow[r,dash] \& e
\end{tikzcd}}
avec $K$ \uline{radiciel} [pas néc. comm!] , on a
\[
H^1(k,G)-H^1(k,G')\quad\text{\uline{surjectif}}.
\]

Par dévissage, on peut supposer $K$ ``de hauteur 1''. Soit $\unsure{\mathfrak{k}|k}$ son Algèbre de
% archive sheet number: 7

\scanpage{9}
Lie, donc $\unsure{G|C}$ opère dans $\unsure{\mathfrak{k}|k}$. Soit $\unsure{\mathfrak{h}|\mathfrak{b}}$ une sous-algèbre de Cartan [existe tjrs], et considérons $N_G(\mathfrak{h})$, alors l'argument déjà utilisé montre que $N_G(\mathfrak{h}).K=G$, donc on \deleted{peut \unsure{remplacer}} a une suite exacte
% slanted
\margin{\uline{grâce au} \uline{thm de} \uline{conjugaison} \uline{pour ss-algèbres} \uline{de Cartan}}
\diagram{\begin{tikzcd}[ampersand replacement=\&]
e \arrow[r,dash] \& N_K(\mathfrak{h}) \arrow[r] \& N_G(\mathfrak{h}) \arrow[r,dash] \& G' \arrow[r,dash] \& e
\end{tikzcd}}
et on peut donc supposer que $G$ \uline{normalise} $\mathfrak{h}$. Mais alors $K$ normalise $\mathfrak{h}$, donc $\mathfrak{k}$ normalise $\mathfrak{h}$, donc $\mathfrak{h}=\mathfrak{k}$, donc $\mathfrak{k}$ est \uline{nilpotente}. Donc $\mathfrak{k}$ se dévisse \deleted{\unsure{en}} par commutateurs, donc (par la \unsure{Théorie} \unsure{Cartier|Cartan}] $K$ aussi se dévisse en groupes commutatifs [N.B. $[\mathfrak{h},\mathfrak{h}]$ est une sous-$p$-algèbre de Lie \unsure{p-fermée} \unsure{?} correspondant au groupe des \unsure{comm.}]. On peut donc supposer $K$ comm., OK.
% archive sheet number: 8

\scanpage{10}
\uline{Lemme 5} Supposons que pour tte \deleted{\unsure{propre}} extension finie \inserted{séparable} $K'$ de $K$, on ait $\operatorname{Br}(K')=0$. Alors on a $H^i(k,\unsure{\mathscr{K}|\mathcal{K}})=0$ pour $i\geq2$, $H$ groupe commutatif fini sur $k$. [N.B. en fait, inutile que $K$ soit fini, comme on voit sans difficulté par dévissage, cf. \deleted{(\unsure{en} \unsure{raison} des \illegible{2} \unsure{groupes} \unsure{lisses} sur $k$)}
% slanted
\margin{Ce raisonnement \unsure{ci-haut} prouve aussi que $H^i(k,U)=0$ pour $i\geq2$ si $U$ unipotent, $k$ corps quelconque. Donc pour $i\unsure{\geq|>}\unsure{2|?}$, \unsure{la seule} cohomologie qui compte est la cohomologie \unsure{étale} \unsure{des|pour} \unsure{groupes} \unsure{algébriques} \unsure{comm.}.}

$\mathscr{K}$ se dévisse en groupes qui sont, soit de t. m. sur $k$, \deleted{soit \unsure{donc} \unsure{à} \unsure{groupes} \unsure{étales}} soit isomorphes à $\alpha_p$. On sait que dans le premier cas, on a une suite exacte
\diagram{\begin{tikzcd}[ampersand replacement=\&]
0 \arrow[r,dash] \& K \arrow[r,dash] \& \unsure{G|C} \arrow[r,dash] \& G' \arrow[r,dash] \& 0
\end{tikzcd}}
avec $G,G'$ des tores, et par Tate--Nakayama on a
\[
H^i(k,G)=H^i(k,G')=0\qquad\text{pour }i>0
\]
[N.B. compte tenu que la cohomologie plate $=$ cohomologie galoisienne]. On en conclut par la suite exacte $H^i(k,K)=0$ pour $i\geq2$. Dans le cas de $\alpha_p$ il est bien connu en tous cas que $H^i(k,\alpha_p)=0$ tt $i$.

\uline{Théorème} Soit $\unsure{K|k}$ un corps tel que $\operatorname{cd}(k)\leq1$. Alors pour tt \deleted{suite exacte $e-K-G-G'$} \inserted{épim. $G-G'$ de} groupes alg sur $k$, $H^1(k,G)\longrightarrow H^1(k,G')$ est un \uline{épimorphisme}
% archive sheet number: 9

\scanpage{11}
\uline{Théorème} \inserted{Soit $k$ un corps \unsure{du|de} \unsure{Thm.}} Soit $K$ un groupe algébrique sur $k$, et $K'$ un noyau sur $k$, localement isomorphe à $K$. Alors $H^2(k,K')=H^2(k,K')'=\deleted{\unsure{\{e\}}}$ uniponctuel.

En effet, $K'$ provient d'un élément de $H^1(k,\underline{\operatorname{Aut}\operatorname{Ext}}(K))$, qui par le raisonnement précédent [\uline{légèrement généralisé}, comme $\underline{\operatorname{Aut}}(K)$ n'est pas néc. un groupe algébrique\unsure{\ldots}] provient d'un élément de $H^1(k,\underline{\operatorname{Aut}}(K))$, ce qui prouve que le noyau $K'$ est représentable par un groupe $K_1$, [loc. isom. à $K$], \deleted{d'où déjà} \inserted{i.e.} \deleted{$\unsure{H^2(k,}$} $H^2(k,K')'\ne\varnothing$. Comme alors $H^2(k,K')\ne\varnothing$ est un fibré principal homogène sous $H^2(k,\unsure{\mathfrak{z}|Z})$, $\mathfrak{z}$ le centre de $\unsure{K_1|K'}$, ou encore de $K$, or on sait que $H^2(k,\mathfrak{z})=0$. Donc $H^2(k,K')$ réduit à un \unsure{pt}, donc $H^2(k,K')'$ aussi [étant non vide et \unsure{dedans contenu}].
% archive sheet number: 10

\scanpage{12}
Noyau $K$ sur le site $X$. Question : $K$ provient-il d'un faisceau en groupes
\[
\text{i.e. }c(K)\in H^2(X,\unsure{I|Z}(K))\text{ est-il dans }\deleted{\unsure{I(K)}}H^2(X,I(K))'\text{ ???}
\]

\deleted{Supposons que l'on puisse trouver un sous-}

À un noyau $K$, est associé le \uline{faisceau} $\unsure{\sigma|\mathscr{T}}(K)$ des \uline{classes} des \uline{sous-groupes de $K$}. \deleted{Supposons \unsure{que nous disposions} d'une section de $\sigma(K)$.}
\inserted{\unsure{On a des opérations} $\sigma(K)\longrightarrow\sigma(K)$, \unsure{telle|t.q.} $\unsure{\varphi|\phi}\longmapsto N(\varphi)$ (\unsure{normalisateur|normaliseur|normalisant}).}
\deleted{Alors on peut \unsure{considérer} le groupe \unsure{normalisateur} dans $\sigma(K)$.}

Chaque fois qu'on a un hom de noyaux $K_0-K$ \deleted{qui est ``injectif''}, il définit une section de $\sigma(K)$ [noyau image].

\deleted{D'ailleurs} Mais une section de $\sigma(K)$ ne définit en général pas un sous-noyau de $K$, sauf celles qui sont stables par l'application $\unsure{\operatorname{Nor}|\operatorname{Na}}$ [``Normalisateur'']. Plus \unsure{généralement|précisément}, si $\varphi$ est une section de $\sigma(K)$, alors $\operatorname{Nor}(\varphi)$ provient d'un sous-noyau de $K$, soit $\underline{\operatorname{Nor}}(\varphi)$, contenant $\unsure{Z|\mathfrak{z}}(K)=\underline{\operatorname{Nor}}(e)$.

\deleted{Supposons que} De plus, on trouve un élément canonique dans $H^2(X,\underline{\operatorname{Nor}}(\varphi))$ dont l'image dans $H^2(X,I(K))$ existe et est $c(K)$ : \deleted{\unsure{il est défini}} c'est la classe de la \unsure{catégorie-loc}
% archive sheet number: 11

\scanpage{13}
2-connexe des faisceaux en groupes, munis d'un isom. de \uline{noyaux} avec $K$, \uline{et} d'un sous-faisceau en groupes dans la classe indiquée par $\unsure{\varphi|\phi}$.

Si $c(K,\varphi)\in H^2(X,\underline{\unsure{\operatorname{Nor}|\operatorname{Na}}}(\varphi))'$, alors $c(K)\in H^2(X,\unsure{I|Z}(K))'$ !

[Principe de \unsure{réduction} de \unsure{Schur|Schreier|schéma}]

Autre principe de réduction. Supposons que nous ayons un hom \inserted{``surjectif''} de noyaux $K\longrightarrow K'$, et que nous ayons donné un isom $K'\simeq\widetilde{F}'$, où $F'$ est un faisceau en groupes.

\deleted{Alors il existe une} Considérons la catégorie locale des faisceaux en groupes $F$, muni d'un \unsure{hom} $F-F'$, et d'un isom $\widetilde{F}\xrightarrow{\sim}K$, \unsure{compatibles|compatible} avec $K\longrightarrow K'$ et $K'\simeq\widetilde{F}'$. C'est une catégorie 2-connexe, et correspond à un noyau $K''$ tel que l'on ait une suite exacte
\diagram{\begin{tikzcd}[ampersand replacement=\&]
e \arrow[r,dash] \& K'' \arrow[r,dash] \& I(K) \arrow[r] \& I(K') \arrow[r,dash] \& e
\end{tikzcd}}
\deleted{Alors} Considérons la classe $c'\in H^2(\unsure{k|X},K'')$. Son image dans \deleted{$\unsure{I(K}$} $H^2(k,I(K))$ est définie et est $c(K)$, donc si $c'\in H^2(X,K'')'$, alors $c\in H^2(X,I(K))'$.
% archive sheet number: 12

\scanpage{14}

$K$ noyau algébrique sur $k$ corps du Thm.

On veut prouver $H^2(k,K)\simeq H^2(k,K)'\simeq\text{un élément}$ [résulte formellement de \uuline{$H^2(k,K)'\ne\varnothing$}, compte tenu que $H^2(k,\unsure{Z|\mathfrak{z}}(K))=0$ !] \deleted{\unsure{Raisonnons sur} \unsure{dimension} de $K$, \unsure{trivial si elle est}}

a) $\dim K=0$. On raisonne par récurrence sur le degré de $K$. Trivial s'il est égal à $1$. S'il est $>1$, par dévissage \inserted{($2^{\text{ième}}$ principe de réduction)}, on est ramené aux cas suivants

(i) $K$ étale et \unsure{on|en} aux deux cas
\[
\left|\begin{aligned}
&K\ \deleted{\text{étale}}\ \text{comm. (trivial)}\\
&[K,K]=K\quad\text{en particulier }K\ \uline{\text{sans centre}}.
\end{aligned}\right.
\]

(ii) $K$ radiciel et $=$ par dévissage
\[
\left\{\begin{aligned}
&K\ \text{radiciel de hauteur }1\\
&\text{et }=\ [K,K]=1\quad\text{en particulier }K\ \text{sans centre}.
\end{aligned}\right.
\]

Le cas (i) correspond à un hom de $\pi$ (groupe de Galois top. de $\bar k/k$] dans $\operatorname{Aut}\operatorname{ext}(A)$, $A$ groupe fini constant. On sait qu'il se remonte à un hom $\pi\unsure{-|\to}\operatorname{Aut}(A)$, [\unsure{Schur|Schreier}+Serre] donc on gagne \unsure{...}

Dans le cas (ii), on considère les sections $\unsure{\varphi|\phi}$ de $\unsure{\sigma|\mathscr{T}}(K)$ correspondant aux ``ss-groupes de Cartan \inserted{infinitésimaux}'' \deleted{qui \unsure{sont leurs}} qui satisfont $\underline{\unsure{\operatorname{Nor}|\operatorname{Na}}}(\varphi)=\varphi$. Par le principe de réduction $1$, on est

% archive sheet number: 13

\scanpage{15}

ramené au cas où $K$ correspond à une $p$-Algèbre de Lie nilpotente, puis, par le principe $2$, au cas où $K$ est commutatif, on gagne encore

b) $\dim K>0$. Par le principe $\unsure{2|1}$, \deleted{\unsure{donc|on}} on est ramené (par les commutateurs) au cas où $K$ est soit commutatif [qui est OK], soit $K=[K,K]$, qui implique $K$ affine. \deleted{\unsure{Enfin,}} Par le même principe [dévissant par un noyau de Frobenius, et utilisant le cas a)] on est ramené au cas où $K$ est simple$/k$.

\deleted{Enfin, \unsure{si on veut}, on peut le \unsure{supposer}}

Prenons alors, la section $\unsure{\varphi|\phi}$ de $\unsure{\sigma|\mathscr{T}}(K)$ correspondant aux tores maximaux, on \deleted{\unsure{donc}} se ramène \inserted{par le principe $1$} au cas où le tore maximal est \deleted{\unsure{donc}} invariant, donc par le principe $2$ aux \deleted{deux} \deleted{\inserted{trois}} \inserted{deux} cas suivants : a) $K$ fini ét., déjà traité b) $K$ \deleted{\unsure{comm.}} unip.\deleted{\unsure{otent}} connexe.

\deleted{Enfin, on peut dévisser par le \unsure{ss. groupe|ss-groupe} \unsure{radiciel|radical}, $=$ tore \unsure{maximal}, et on est \unsure{ramené}.} Alors on est ramené par dévissage par commutateurs, au cas comm., qui est OK.

% archive sheet number: 14

\scanpage{16}

\deleted{\unsure{Égalité} \illegible{2}}

Cas non $\frac12$ simple : \unsure{Questions|Question} de rationalité des \unsure{radic\textsuperscript{l}|radic\textsuperscript{t}}, définition de la ``variété des drapeaux'' etc.

% archive sheet number: 15

\scanpage{17}

\unsure{Soient|Soit} $a\in k$, et $G$ le groupe
\[
G=\left\{\left(\begin{array}{cc}
\begin{matrix}x&ay\\y&x\end{matrix}&0\\
0&\begin{matrix}u&v\\w&z\end{matrix}
\end{array}\right)\ \middle|\ 
\begin{vmatrix}x&ay\\y&x\end{vmatrix}=\begin{vmatrix}u&v\\w&z\end{vmatrix}\ne0\right\}.
\]
\[
\left(\begin{array}{cc}
\begin{matrix}x&ay\\y&x\end{matrix}&0\\
0&\begin{matrix}u&v\\w&z\end{matrix}
\end{array}\right)\rightsquigarrow
\begin{pmatrix}u&v\\w&z\end{pmatrix}/(\text{centre})\quad\text{est un hom.}
\]
sur $GP(2)$ dont le noyau est
\[
\begin{aligned}
Z&=\left\{\left(\begin{array}{cc}
\begin{matrix}x&ay\\y&x\end{matrix}&0\\
0&\begin{matrix}u&0\\0&u\end{matrix}
\end{array}\right)\ \middle|\ x^2-ay^2=u^2\ne0\right\}\\
&=\text{centre de }G=\text{radical de }G.
\end{aligned}
\]

$G$ est déf. sur $k$, \unsure{est|et} si $\text{caract.}=2$, $a\notin k^2$, $Z$ ne l'est pas.

$\exists$ exemples analogues pour chaque caract. $p\ne0$.

% archive sheet number: 16

